% Options for packages loaded elsewhere
\PassOptionsToPackage{unicode}{hyperref}
\PassOptionsToPackage{hyphens}{url}
%
\documentclass[
]{article}
\usepackage{amsmath,amssymb}
\usepackage{lmodern}
\usepackage{iftex}
\ifPDFTeX
  \usepackage[T1]{fontenc}
  \usepackage[utf8]{inputenc}
  \usepackage{textcomp} % provide euro and other symbols
\else % if luatex or xetex
  \usepackage{unicode-math}
  \defaultfontfeatures{Scale=MatchLowercase}
  \defaultfontfeatures[\rmfamily]{Ligatures=TeX,Scale=1}
\fi
% Use upquote if available, for straight quotes in verbatim environments
\IfFileExists{upquote.sty}{\usepackage{upquote}}{}
\IfFileExists{microtype.sty}{% use microtype if available
  \usepackage[]{microtype}
  \UseMicrotypeSet[protrusion]{basicmath} % disable protrusion for tt fonts
}{}
\makeatletter
\@ifundefined{KOMAClassName}{% if non-KOMA class
  \IfFileExists{parskip.sty}{%
    \usepackage{parskip}
  }{% else
    \setlength{\parindent}{0pt}
    \setlength{\parskip}{6pt plus 2pt minus 1pt}}
}{% if KOMA class
  \KOMAoptions{parskip=half}}
\makeatother
\usepackage{xcolor}
\setlength{\emergencystretch}{3em} % prevent overfull lines
\providecommand{\tightlist}{%
  \setlength{\itemsep}{0pt}\setlength{\parskip}{0pt}}
\setcounter{secnumdepth}{-\maxdimen} % remove section numbering
\ifLuaTeX
  \usepackage{selnolig}  % disable illegal ligatures
\fi
\IfFileExists{bookmark.sty}{\usepackage{bookmark}}{\usepackage{hyperref}}
\IfFileExists{xurl.sty}{\usepackage{xurl}}{} % add URL line breaks if available
\urlstyle{same} % disable monospaced font for URLs
\hypersetup{
  hidelinks,
  pdfcreator={LaTeX via pandoc}}

\author{Dietmar Gerald Schrausser}
\date{09.05.2026}


\begin{document}

\hypertarget{folium-irfo.-.1.r}{%
\section{Foliu{[}m{]} Ir/Fo. .1.r}\label{folium-irfo.-.1.r}}

A \emph{contemporary} English translation is presented from comparing
the existing Latin, Early New High German (ENHG) and English
(translation) texts for better comprehensibility. Especially since the
ENHG text is difficult to understand compared to modern High German, as
a \emph{profound} knowledge of German (as a mother tongue) as well as
already aknowledge of various \emph{dialects} is required.

\hypertarget{foreword-a-brief-description-of-the-work-of-the-six-days-of-the-creation-of-the-world.}{%
\subsection{\texorpdfstring{Foreword: A brief description of the
\emph{work of the six days} of the creation of the
world.}{Foreword: A brief description of the work of the six days of the creation of the world.}}\label{foreword-a-brief-description-of-the-work-of-the-six-days-of-the-creation-of-the-world.}}

Since\footnote{`\textbf{C}Um' and `\textbf{D}Ieweil', c.f. Diodorus
  (\href{https://doi.org/10.3931/e-rara-35702}{1578}, Book I, pp.5-6).}
among the most learned and capable\footnote{`prestantissimos' and
  `fürnamste{[}n{]}'.} men who have described the true nature and
history of (i) the origin of the world, as well as the (ii) first
generation\footnote{`generatione' and ' Geburt'.} of humans, there is a
\emph{divided} opinion, we would like to report on this very briefly,
beginning with the \emph{primeval times}\footnote{`priscis' and
  `vordern'.} and insofar as this seems appropriate\footnote{`licebit'
  and `gezime{[}n{]} wil'.} for such long-past\footnote{`remotis' and
  `enthlegne{[}n{]}'.} things.

The view that the world (i) was not created and is indestructible, that
humanity has existed since eternity and has no origin, contrasts with
the assertion that the world (ii) was created and is transient, and that
humankind has possessed it from the first generation. It was the highly
educated Greeks who collected all historical records and followed the
theory that before the beginning of all things, heaven and earth, while
they were still united, existed in only a single form. Only afterward,
through the separation and division of matter, did the world assume the
order and structure in which we see it today.

\begin{enumerate}
\def\labelenumi{(\Roman{enumi})}
\tightlist
\item
  These (Greeks) claim that (i) the air and the fiery part of the upper
  layers were constantly in motion, becoming lighter, and that the sun
  and the many stars were supported by them. The (ii) dark and solid
  part, however, along with the moist substances, was carried by its own
  weight into the lowest layer. After these substances had mixed, the
  sea arose from the mist, and the solid, clayey, and soft matter became
  earth. As the earth became denser from the heat of the sun, rotting
  mud formed, covered by a thin skin. From such swamps and pools,
  diverse life forms arose. Those that had absorbed more heat became
  winged creatures and rose into the upper layers, while the drier and
  heavier ones became crawling, terrestrial animals. Those that assumed
  a water-bound nature were carried into the element that was
  appropriate for their kind.
\end{enumerate}

As the earth dried out from the heat of the sun and the action of the
air, a mixture of more perfect creatures, male and female, arose. This
is recounted by \textbf{Euripides}, the tragedian and \emph{student} of
\textbf{Anaxagoras}, the master of natural history. Likewise, humans are
said to have been born in open fields, wandering far and wide, leading a
wild and disorderly life, nourished by herbs and the fruits of
trees.\footnote{Beginning of the first book of Diodorus Siculus'
  `Bibliotheca historica', 1st c.~BC (c.f. Diodorus,
  \href{https://doi.org/10.3931/e-rara-68988}{1476}, fol.~15v).}

However, since much (both old and new) has been written on this subject
in Latin, Greek, Chaldean, and Hebrew, we wish to leave behind these
old, \emph{erroneous} accounts and turn to the \emph{profound} Mosaic
writings concerning the creation of the world and the work of the six
days, in which the mysteries of all nature are encompassed. For
\textbf{Moses}, the prophet and father of God's historians, understood
all this through the inspiration of the Holy Spirit, the master of all
truth. His human understanding and experience in all areas of knowledge
have been attested to by our people, his people, and also by the
Gentiles.

King \textbf{Solomon}, an interpreter of nature and living beings,
reveals in his Book of Wisdom that he derived his knowledge of these
essential phenomena from the Mosaic laws. He was (as \textbf{Luke} and
\textbf{Philo}, our most renown teachers, report) thoroughly familiar
with the entire Egyptian tradition. According to \textbf{Hermippus},
\textbf{Pythagoras} also drew much of his philosophy from Mosaic law.

The philosopher \textbf{Numenius} claims that \textbf{Plato} was a true
Attic Moses, since a treasure of true wisdom lies hidden in the
beginnings of his natural philosophy. He speaks learnedly and wisely
about all things as originating from God, about their relationships to
one another, their number and order, as well as their transformations.
It was therefore considered a law among the ancient Hebrews (as
\textbf{Jerome} also holds) that one should not concern oneself with
this creation story before reaching a mature age\footnote{`nisi matura
  iam etate atingeret' and `niemant dan der zeitigs alters wer', older
  than 15 years.}. But \emph{what} the most pious men, \textbf{Ambrose}
and \textbf{Augustine}, \textbf{Strabo}, also \textbf{Bede} or
\textbf{Remigius} and the younger ones, \textbf{Aegidius},
\textbf{Albertus}, as well as the Greeks \textbf{Philo},
\textbf{Origen}, \textbf{Basil}, \textbf{Theodorus},
\textbf{Apollinarius}, \textbf{Didymus}, \textbf{Gennadius},
\textbf{Chrysostome} etc., have written about this book, we will leave
untouched.

Nor will we mention what \textbf{Jonetes}, \textbf{Anchelos}, or
\textbf{Simeon the Elder} said in Chaldean, or what the Hebrews
\textbf{Eleazadus}, \textbf{Aba}, \textbf{John}, \textbf{Neonius},
\textbf{Isaac}, \textbf{Josephus}, \textbf{Gersonides}, \textbf{Sadias},
\textbf{Abraham}, etc., wrote. Rather, we will briefly outline the order
of the six days according to Moses, in which the \emph{divine} creation
of the world took place, as recounted in the sacred and profound
religious scriptures.

\begin{enumerate}
\def\labelenumi{(\Roman{enumi})}
\setcounter{enumi}{1}
\tightlist
\item
  When \emph{God} created the world, he placed his first and greatest
  son at the head of his infinite work and appointed him as advisor and
  master craftsman in the planning, beautification and design of things,
  since he was endowed with wisdom and understanding in abundance.
\end{enumerate}

However, the most pertinent question is, from what did God create all
this great and wonderful thing, \emph{seemingly} out of nothing. At this
point, it is necessary to disregard transient things and turn our gaze
to the \emph{seat}, to where the \emph{true} God dwells. \emph{He} gave
the earth eternal stability, hung shining stars in the sky, created the
brightest suns, surrounded the earth with the sea, made the rivers flow,
the fields spread, the valleys descend, the forests adorn themselves
with leaves, and the rocky mountains rise. All this was created not by
\textbf{Jupiter}\footnote{\begin{enumerate}
  \def\labelenumi{\alph{enumi}.}
  \setcounter{enumi}{18}
  \tightlist
  \item
    Boccaccio
    (\href{https://digi.vatlib.it/view/MSS_Pal.lat.938}{1360}).
  \end{enumerate}}, but by the master craftsman of the world, the
\emph{source} of all that is best, that is called \emph{God} and whose
beginning is incomprehensible and unfathomable.

The ancients spoke of three \emph{worlds}: (i) a supreme one, the world
of angels; (ii) a heavenly world; and (iii) a world beneath the moon, in
which we live; a world of \emph{darkness}, which is nevertheless
regularly illuminated by the light of heaven. Beyond this, there is a
fourth world in which all the qualities of the other worlds are united:
(iv) humankind itself. In school, we learned that humankind is a
microcosm in which the likeness of God can be recognized through (i) the
body composed of the elements and (ii) a heavenly spirit, through (iii)
the growth soul of plants, (iv) the sentience of rational animals and
through (v) reason and the angelic\footnote{intellectual.} mind.

Moses gives a detailed account of the \emph{three} worlds and how God
ordered them. On the mountain, where he learned, he (Moses) was also
commanded in detail (as we can read it) to do all things according to
the parable he saw there. We will now present an overview of \emph{what}
the Book of Moses teaches about the completed works of the six
days.\footnote{from the `Heptaplus' (Mirandola \& Mirandola,
  \href{https://doi.org/10.3931/e-rara-61217}{1601}, p.~1 ff.).}

\hypertarget{references}{%
\subsection{References}\label{references}}

Boccaccio, G. (1360). \emph{Genealogiae deorum gentilium libri XV}. Pal.
lat. 938. Vatican: Biblioteca Apostolica Vaticana.
\url{https://digi.vatlib.it/view/MSS_Pal.lat.938}

Diodorus. (1476). \emph{Diodori Sicvli Historiarvm Priscarvm}. Edited by
Poggio Bracciolini, G. F., \& Tacitus, C. Venetiis: per Andreā Iacobi
Katharēsem Andrea Vendramino Duce fortunatissimo.
\url{https://doi.org/10.3931/e-rara-68988}

---------. (1578). \emph{Diodori Siculi Bibliothecae Historicae Libri
XV: \ldots{} Nunc Iterum Latinè \ldots{} Recogniti, \& \ldots{}
Eduntur}. Edited by Grynaeus, J. J., Châteillon, J., Dictys, Dares, \&
Tryphiodorus. Basileae: ex officina Henricpetrina.
\url{https://doi.org/10.3931/e-rara-35702}

Mirandola, G., \& Mirandola, G. F. (1601). \emph{Ioannis Pici,
Mirandulae \ldots{} opera quae extant omnia \ldots{}} Basileae: per
Sebastianum Henricpetri. \url{https://doi.org/10.3931/e-rara-61217}

\end{document}
